\begin{table}[H]
\centering
\caption{Standardized Effect Sizes for Main Outcomes}
\label{tab:sde}
\begin{threeparttable}
\footnotesize
\begin{tabular}{@{}lrrcrrcl@{}}
\toprule
Outcome & $\hat{\beta}$ & SE & SD($X$) & SD($Y$) & SDE & SE(SDE) & Classification \\
\midrule
Separations (Low-Wage) & 0.0026 & 0.0180 & --- & 0.9554 & 0.0027 & 0.0188 & Null \\
Separations (High-Wage) & -0.0346 & 0.0224 & --- & 1.2204 & -0.0283 & 0.0184 & Small negative \\
Employment (Low-Wage) & -0.0161 & 0.0131 & --- & 1.0315 & -0.0156 & 0.0127 & Small negative \\
Separations (DDD) & -0.0050 & 0.0541 & --- & 1.3372 & -0.0037 & 0.0405 & Null \\
\bottomrule
\end{tabular}
\begin{tablenotes}[flushleft]
\scriptsize
\item \textit{Notes:} \textbf{Country:} United States. \textbf{Research question:} Do state increases in the UI taxable wage base reduce employer-initiated separations, particularly in low-wage industries where the marginal tax cost increase is binding? \textbf{Policy mechanism:} State UI systems are experience-rated: employers who lay off more workers face higher payroll tax rates. Raising the taxable wage base increases the per-worker tax liability for each separation, making layoffs more expensive---especially for workers earning below the new threshold. \textbf{Outcome definition:} Log quarterly separations and log quarterly employment from Census Quarterly Workforce Indicators (QWI) at the state-industry-quarter level. \textbf{Treatment:} Binary indicator for state having raised its UI taxable wage base above the \$7,000 federal minimum by more than \$3,000. \textbf{Data:} Census QWI, 2001Q1--2023Q4, state-industry-quarter level. \textbf{Method:} Two-way fixed effects DiD and triple-difference with state, quarter, and industry fixed effects; state-clustered standard errors. \textbf{Sample:} States that raised wage bases vs.\ states at federal minimum; low-wage industries (retail, food) vs.\ high-wage industries (finance, professional). SDE $= \hat{\beta} / \text{SD}(Y)$ where SD($Y$) is pre-treatment SD. Classification refers to magnitude, not statistical significance: Large ($|$SDE$| > 0.15$), Moderate ($0.05$--$0.15$), Small ($0.005$--$0.05$), Null ($< 0.005$).
\end{tablenotes}
\end{threeparttable}
\end{table}

